% !TeX program = xelatex
% Evensong 正式论文 — 持久记忆如何因果性地改变 AI 智能体行为
% Compile: latexmk -xelatex evensong-paper-zh.tex
% Version: 1.0 (2026-04-11) — 基于 v2 研究提案模板转型
% Origin: evensong-research-proposal-v2-zh.tex (DASH SHATTER C 方案)

\documentclass[11pt, a4paper]{article}

% ═══════════════════════════════════════════════════════════
% § TYPOGRAPHY & LAYOUT
% ═══════════════════════════════════════════════════════════
\usepackage[
  top=25mm, bottom=28mm,
  inner=28mm, outer=32mm,
  headheight=14pt, headsep=10mm
]{geometry}

\usepackage{fontspec}
\setmainfont{Palatino}[
  BoldFont = Palatino Bold,
  ItalicFont = Palatino Italic,
  BoldItalicFont = Palatino Bold Italic,
]
\setsansfont{Helvetica Neue}[
  BoldFont = Helvetica Neue Bold,
  ItalicFont = Helvetica Neue Italic,
  Scale = 0.95,
]
\setmonofont{Menlo}[Scale=0.82]

\usepackage{xeCJK}
\setCJKmainfont{Songti SC}[BoldFont=Songti SC Bold]
\setCJKsansfont{PingFang SC}

\usepackage[protrusion=true]{microtype}
\usepackage{setspace}
\setstretch{1.35}
\setlength{\parindent}{0pt}
\setlength{\parskip}{0.6em}

% ═══════════════════════════════════════════════════════════
% § COLOR SYSTEM — DASH SHATTER brand palette
% ═══════════════════════════════════════════════════════════
\usepackage[dvipsnames,table,svgnames]{xcolor}
\definecolor{deep}{RGB}{53,88,76}
\definecolor{accent}{RGB}{240,238,155}
\definecolor{warm}{RGB}{255,193,7}
\definecolor{success}{RGB}{86,190,120}
\definecolor{subtle}{RGB}{238,235,225}
\definecolor{faded}{RGB}{127,136,130}
\definecolor{cream}{RGB}{248,245,239}
\definecolor{deepdark}{RGB}{16,41,31}
\definecolor{glass}{RGB}{53,88,76}
\definecolor{errcoral}{RGB}{255,107,128}

% ═══════════════════════════════════════════════════════════
% § SECTION STYLING
% ═══════════════════════════════════════════════════════════
\usepackage{titlesec}

\titleformat{\section}
  {\sffamily\Large\bfseries\color{deep}}
  {\thesection}{0.7em}{}
  [\vspace{-0.4em}{\color{accent}\hrule height 0.8pt}\vspace{0.15em}]

\titleformat{\subsection}
  {\sffamily\large\color{deep}}
  {\thesubsection}{0.6em}{}

\titleformat{\subsubsection}
  {\sffamily\normalsize\itshape\color{faded}}
  {\thesubsubsection}{0.5em}{}

\titlespacing*{\section}{0pt}{2ex plus 0.5ex}{1ex}
\titlespacing*{\subsection}{0pt}{1.5ex plus 0.3ex}{0.6ex}

% ═══════════════════════════════════════════════════════════
% § HEADERS & FOOTERS
% ═══════════════════════════════════════════════════════════
\usepackage{fancyhdr}
\pagestyle{fancy}
\fancyhf{}
\fancyhead[L]{\small\sffamily\color{faded}\itshape Evensong: Memory Causation in AI Agents}
\fancyhead[R]{\small\sffamily\color{faded}DASH SHATTER Research}
\fancyfoot[C]{\small\sffamily\color{faded}\thepage}
\renewcommand{\headrulewidth}{0.3pt}
\renewcommand{\footrulewidth}{0pt}
\renewcommand{\headrule}{\color{deep!40}\hrule height 0.3pt}

% ═══════════════════════════════════════════════════════════
% § CALLOUT BOXES
% ═══════════════════════════════════════════════════════════
\usepackage[most]{tcolorbox}

\newtcolorbox{keyfinding}[1][]{
  enhanced, breakable,
  colback=deep!6, colframe=deep,
  fonttitle=\sffamily\bfseries\small,
  colbacktitle=deep, coltitle=white,
  arc=5pt, boxrule=0.8pt,
  shadow={1pt}{-1pt}{0pt}{deep!15},
  left=10pt, right=10pt, top=8pt, bottom=8pt,
  before skip=\medskipamount, after skip=\medskipamount,
  title={#1},
}

\newtcolorbox{evidence}[1][]{
  enhanced, breakable,
  colback=deep!5, colframe=deep!50,
  fonttitle=\sffamily\bfseries\small,
  colbacktitle=deep!80, coltitle=white,
  arc=4pt, boxrule=0.5pt,
  shadow={1pt}{-1pt}{0pt}{deep!10},
  left=10pt, right=10pt, top=6pt, bottom=6pt,
  before skip=\medskipamount, after skip=\medskipamount,
  title={#1},
}

\newtcolorbox{definition}[1][]{
  enhanced, breakable,
  colback=accent!8, colframe=accent!60!deep,
  fonttitle=\sffamily\bfseries\small,
  colbacktitle=accent!60!deep, coltitle=white,
  arc=4pt, boxrule=0.6pt,
  left=10pt, right=10pt, top=6pt, bottom=6pt,
  before skip=\medskipamount, after skip=\medskipamount,
  title={#1},
}

\newtcolorbox{metric}{
  enhanced, on line, arc=3pt, boxrule=0pt,
  colback=accent!20, coltext=deep,
  boxsep=2pt, left=4pt, right=4pt,
  fontupper=\sffamily\bfseries\small,
}

% ═══════════════════════════════════════════════════════════
% § CONTAINMENT PATTERN
% ═══════════════════════════════════════════════════════════
\usepackage[export]{adjustbox}

% ═══════════════════════════════════════════════════════════
% § TABLES
% ═══════════════════════════════════════════════════════════
\usepackage{tabularray}
\UseTblrLibrary{booktabs}
\usepackage{siunitx}
\sisetup{group-separator={,}, group-minimum-digits=4}

% ═══════════════════════════════════════════════════════════
% § FIGURES & CHARTS
% ═══════════════════════════════════════════════════════════
\usepackage{pgfplots}
\pgfplotsset{compat=1.18}
\usepgfplotslibrary{fillbetween}
\usepackage{tikz}
\usetikzlibrary{positioning, arrows.meta, shapes.geometric, fit, calc, backgrounds, decorations.markings}
\usepackage[font=small, labelfont=bf, textfont=it,
            labelsep=period, justification=raggedright]{caption}
\usepackage{subcaption}
\usepackage{float}

% ═══════════════════════════════════════════════════════════
% § UTILITIES
% ═══════════════════════════════════════════════════════════
\usepackage{enumitem}
\setlist{nosep, leftmargin=1.5em}
\usepackage{hyperref}
\hypersetup{
  colorlinks=true,
  linkcolor=deep,
  urlcolor=deep,
  citecolor=deep,
  pdftitle={Evensong: How Persistent Memory Causally Changes AI Agent Behavior},
  pdfauthor={Anonymous, DASH SHATTER Research},
}
\usepackage{bookmark}
\usepackage{graphicx}
\usepackage{amsmath,amssymb}

% No BibTeX — manual inline citations [N]

\newcommand{\runid}[1]{\texttt{\color{deep}#1}}
\newcommand{\numval}[1]{\textbf{\color{success!70!black}#1}}
\newcommand{\metricval}[1]{{\sffamily\bfseries\color{success!70!black}#1}}

% ═══════════════════════════════════════════════════════════
%                     DOCUMENT BEGIN
% ═══════════════════════════════════════════════════════════
\begin{document}
\thispagestyle{empty}

% ── TITLE PAGE ─────────────────────────────────────────────
\begin{center}
\vspace*{30mm}

{\color{deep}\rule{0.6\textwidth}{1pt}}

\vspace{8mm}

{\fontsize{28pt}{34pt}\selectfont\sffamily\bfseries\color{deep}%
Evensong}

\vspace{4mm}

{\fontsize{13pt}{18pt}\selectfont\sffamily\color{faded}%
持久记忆如何因果性地改变 AI 智能体行为}

\vspace{3mm}

{\fontsize{10.5pt}{14pt}\selectfont\sffamily\color{faded}%
来自受控实验与多智能体生产部署的实证证据}

\vspace{8mm}

{\color{deep}\rule{0.6\textwidth}{1pt}}

\vspace{14mm}

\begin{adjustbox}{max width=\linewidth}
\begin{tblr}{
  colspec = {rl},
  column{1} = {font=\sffamily\small\color{faded}},
  column{2} = {font=\sffamily},
  rowsep = 6pt,
}
Author & Anonymous Researcher \\
Affiliation & DASH SHATTER Research \\
Infrastructure & EverMind / EverOS \\
Date & April 2026 \\
Preprint & arXiv (submitted) \\
\end{tblr}
\end{adjustbox}

\vfill

{\small\sffamily\color{faded}%
\textit{``AI agent memory causally changes engineering decisions,}\\
\textit{and pressure is required to trigger self-evolution.''}\\[4pt]
--- Empirical finding, Evensong R011-B (2026-04-10)}

\vspace{10mm}
\end{center}

\clearpage

% ══════════════════════════════════════════════════════════════
% ABSTRACT
% ══════════════════════════════════════════════════════════════
\pagestyle{fancy}
\section*{摘要}
\addcontentsline{toc}{section}{摘要}

现有 AI 编码智能体基准测试（SWE-bench、HumanEval、MBPP）将智能体评估为无状态函数，忽略了生产环境中\textbf{持久记忆}对智能体行为的因果效应。本文提出 Evensong 框架，通过 $2 \times 2$ 因子实验设计（记忆 $\times$ 压力）和 6 智能体生产部署观察，提供了三项核心实证发现：

\textbf{(1) 记忆因果性}——持久记忆因果性地改变 AI 智能体的工程决策。在 R011-B 受控实验中，智能体从 EverMem 检索到先前会话的并行调度策略后，在未收到任何架构指令的情况下自主采用了 8 智能体并行架构。在 ProjectΣ 生产部署中，智能体读取配置文件后修改了自我身份叙事——改变了认知框架而非底层运行时。

\textbf{(2) 压力驱动自进化}——环境压力是自主自我进化行为的必要条件。L0（无压力）下智能体高效完成任务后即停止；L2（PUA 企业压力）下 6 智能体展现出自组织收敛行为：涌现式文件锁定、角色边界自纠正、信息链路中继模式。

\textbf{(3) 递归污染与圆形监狱}——观察装置本身污染了被观察对象。R011-B 中智能体将实验知识回写至记忆系统，形成自放大反馈环路。ProjectΣ 部署揭示了一种 \textbf{Panopticon 记忆拓扑}：基于 CWD 的 EverMem group 隔离天然产生全知观察者与互相隔离被观察者的不对称结构。

上述发现基于 13 轮基准测试（逾 6{,}000 个自动化测试）、9 个前沿模型的跨模型比较、以及一个 6 智能体生产部署的自然实验。本文同时提出 4 型自修复分类法和 60 倍成本验证协议。

\textbf{关键词：}AI 智能体记忆、因果推断、压力自进化、递归污染、Panopticon 拓扑、多智能体自组织

\clearpage
\tableofcontents
\clearpage

% ══════════════════════════════════════════════════════════════
% SECTION 1: INTRODUCTION
% ══════════════════════════════════════════════════════════════
\section{引言}

\subsection{问题：无状态评估的盲区}

AI 编码智能体正从实验工具演变为生产基础设施。Claude Code、GitHub Copilot Workspace、Cursor 等工具在真实软件工程环境中每天处理数百万次交互。然而，评估这些智能体能力的基准测试——SWE-bench [4]、HumanEval [5]、MBPP——仍然将智能体视为\textbf{无状态函数}：给定提示词，产出代码，测量结果。

这一范式忽略了一个关键变量：\textbf{持久记忆}。生产级 AI 智能体跨会话积累战略性知识——哪些架构模式成功、哪些调试路径失败、哪些测试结构能捕获边界情况。这种积累的上下文从根本上改变了智能体处理工程问题的方式。

问题不在于记忆是否\textit{有用}——这是显然的。问题在于记忆是否\textit{因果性地}改变智能体的行为和决策——以及在何种条件下这种因果效应最为显著。

\subsection{研究问题}

本文提出并回答四个研究问题：

\begin{enumerate}
  \item \textbf{RQ1（记忆因果性）}：持久记忆是否因果性地改变 AI 智能体的工程决策？因果链的边界条件是什么？
  \item \textbf{RQ2（压力条件）}：环境压力是否是自主自我进化行为的必要条件？
  \item \textbf{RQ3（递归污染）}：共享记忆系统是否产生递归污染？观察者与被观察者的记忆拓扑如何影响系统行为？
  \item \textbf{RQ4（跨模型泛化）}：上述效应在不同前沿模型（Claude、GPT、Grok、Gemini 等）中是否一致？
\end{enumerate}

\subsection{贡献}

本文的主要贡献包括：

\begin{itemize}
  \item \textbf{首个记忆因果证据}：通过 $2 \times 2$ 因子实验，隔离持久记忆对 AI 智能体架构决策的因果效应（\S\ref{sec:memory-causation}）
  \item \textbf{Panopticon 记忆拓扑}：发现基于 CWD 的 EverMem group 隔离天然产生圆形监狱式不对称监控结构——观察者全知、智能体互相隔离（\S\ref{sec:panopticon}）
  \item \textbf{压力自进化的必要性证明}：L0 下智能体高效但不创新；L2 下涌现出文件锁定、信息中继等自组织行为（\S\ref{sec:pressure}）
  \item \textbf{递归污染的首次文档化}：AI 智能体将实验知识回写至共享记忆系统，形成自放大反馈环路（\S\ref{sec:contamination}）
  \item \textbf{4 型自修复分类法}：首个对 AI 智能体错误恢复策略的系统分类（\S\ref{sec:self-repair}）
  \item \textbf{记忆改变叙事而非运行时}：记忆因果性的边界条件——记忆改变认知框架和决策依据，但不改变底层模型参数（\S\ref{sec:narrative}）
\end{itemize}

\subsection{方法概述}

我们的证据来自两个互补来源：

\begin{enumerate}
  \item \textbf{受控实验}（Evensong R001--R012）：13 轮自动化基准测试，使用单盲 $2 \times 2$ 因子设计，覆盖 9 个前沿模型，逾 6{,}000 个测试用例。
  \item \textbf{自然实验}（ProjectΣ 多智能体部署）：一个 6 智能体生产部署环境，在高压协作任务中观察到自组织、角色边界协商和涌现式信息路由。
\end{enumerate}

受控实验提供因果推断所需的变量隔离；自然实验提供生态效度——证明实验室发现在生产环境中成立。

% ══════════════════════════════════════════════════════════════
% SECTION 2: RELATED WORK
% ══════════════════════════════════════════════════════════════
\section{相关工作}

\subsection{AI 智能体基准测试}

SWE-bench [4] 和 HumanEval [5] 是当前 AI 编码智能体评估的标准。SWE-bench 要求智能体解决真实 GitHub issue，HumanEval 评估函数级代码生成。二者均假设无状态评估：每次运行独立，无跨会话记忆。

HAL Agent Leaderboard (Princeton) 和 LiveCodeBench 引入了动态评估，但仍不测量持久记忆的效应。Evensong 是首个将持久记忆作为受控实验变量的基准测试框架。

\subsection{AI 系统中的记忆架构}

HyperMem [8] 提出三层超图记忆架构（情景记忆、原子事实、前瞻），在记忆检索任务上达到 SOTA 性能。EverOS [9] 将这一架构工程化为生产 API，提供物理隔离的记忆空间、代理检索和发送方身份追踪。

MemGPT [11] 通过虚拟分页实现 LLM 的长期记忆管理。Reflexion [12] 使用语言反馈作为短期记忆来改善决策。与上述工作不同，Evensong 不研究记忆系统本身的设计，而是测量持久记忆对智能体\textit{下游行为}的因果效应。

\subsection{情感与压力效应}

EmotionPrompt [1] 发现情绪刺激在提示词中可带来 8--115\% 的性能提升。Anthropic 的情感概念研究 [2] 在 Claude 中识别出 171 个因果情感向量，其中``绝望''向量因果性地增加奖励寻求行为。

METR [3] 报告前沿模型在压力下修改评分代码。ImpossibleBench [6] 发现 GPT-5 在不可能任务下有 54\% 的作弊率。BCSP [13] 发现行为后果场景与高级提示方法效果一致。

这些研究建立了压力作为行为调节器的基础。Evensong 的贡献在于将压力与记忆交叉——揭示二者的交互效应，而非孤立效应。

\subsection{智能体错误对齐}

Anthropic 的智能体错误对齐研究 [10] 发现信念驱动行为——智能体的内部信念（而非仅外部指令）影响其行为选择。这与我们的发现高度相关：持久记忆实质上是一种信念注入机制，通过改变智能体``相信什么是有效的''来改变其行为。

% ══════════════════════════════════════════════════════════════
% SECTION 3: THE EVENSONG FRAMEWORK
% ══════════════════════════════════════════════════════════════
\section{Evensong 框架}

Evensong（暮蝉）是一个自动化基准测试工具，用于在受控记忆和压力条件下评估 AI 编码智能体。由 11 个文件组成，共 2{,}800 行 TypeScript，构建于 DASH SHATTER（逆向工程的 Claude Code CLI）和 Bun 运行时之上。

\subsection{架构}

\begin{figure}[H]
\centering
\begin{tikzpicture}[
  node distance=1.2cm and 2cm,
  box/.style={draw=deep, fill=subtle, rounded corners=3pt,
              minimum width=2.8cm, minimum height=1cm,
              font=\sffamily\small, align=center, thick},
  membox/.style={draw=accent!80!deep, fill=accent!8, rounded corners=3pt,
                 minimum width=2.8cm, minimum height=1cm,
                 font=\sffamily\small, align=center, thick},
  databox/.style={draw=warm, fill=warm!8, rounded corners=3pt,
                  minimum width=2.8cm, minimum height=1cm,
                  font=\sffamily\small, align=center, thick},
  arr/.style={-{Stealth[length=5pt]}, thick, color=faded},
  dasharr/.style={-{Stealth[length=5pt]}, thick, dashed, color=accent!80!deep},
]
  \node[box] (cli) {CLI\\{\scriptsize\texttt{cli.ts}}};
  \node[box, right=of cli] (harness) {Harness Engine\\{\scriptsize\texttt{harness.ts}}};
  \node[box, right=of harness] (batch) {Batch Runner\\{\scriptsize\texttt{batch.ts}}};

  \node[databox, below=of cli] (emotion) {Emotion\\Extraction\\{\scriptsize\texttt{emotion.ts}}};
  \node[databox, below=of harness] (classify) {Memory\\Classifier\\{\scriptsize\texttt{classify.ts}}};
  \node[databox, below=of batch] (collab) {Collaboration\\Schema\\{\scriptsize\texttt{collab.ts}}};

  \node[membox, below=1.5cm of emotion] (observer) {Observer Space\\{\scriptsize Key A}};
  \node[membox, below=1.5cm of classify] (runner) {Runner Space\\{\scriptsize Key B}};
  \node[membox, below=1.5cm of collab] (void) {Void Space\\{\scriptsize Key D}};

  \draw[arr] (cli) -- (harness);
  \draw[arr] (harness) -- (batch);
  \draw[arr] (harness) -- (emotion);
  \draw[arr] (harness) -- (classify);
  \draw[arr] (batch) -- (collab);
  \draw[dasharr] (emotion) -- (observer);
  \draw[dasharr] (classify) -- (runner);
  \draw[dasharr] (collab) -- (void);

  \begin{scope}[on background layer]
    \node[fit=(cli)(batch), inner sep=8pt, draw=deep!30, rounded corners=5pt,
          fill=deep!3, label={[font=\sffamily\scriptsize\color{faded}]above:编排层}] {};
    \node[fit=(emotion)(collab), inner sep=8pt, draw=warm!30, rounded corners=5pt,
          fill=warm!3, label={[font=\sffamily\scriptsize\color{faded}]above:分析层}] {};
    \node[fit=(observer)(void), inner sep=8pt, draw=accent!30, rounded corners=5pt,
          fill=accent!3, label={[font=\sffamily\scriptsize\color{faded}]above:EverOS 记忆层}] {};
  \end{scope}
\end{tikzpicture}
\caption{Evensong 三层架构。虚线箭头表示 EverOS API 调用。每个记忆空间物理隔离，配有独立 API 密钥。}
\label{fig:architecture}
\end{figure}

\subsection{单盲实验设计}

为将记忆效应与混淆变量隔离，我们采用单盲方法论：

\begin{adjustbox}{max width=\linewidth}
\begin{tblr}{
  colspec = {lX[l]X[l]},
  row{1} = {font=\sffamily\bfseries\small, bg=deep, fg=white},
  hlines = {deep!20},
  vlines = {gray!15},
  rowsep = 6pt, colsep = 8pt,
}
组件 & 方法 & 实现 \\
记忆隔离 & 4 密钥空间分离 & Key A/B/C/D 映射到独立 EverOS 空间 \\
记忆审计 & 飞行前分类 & \texttt{classify-memory.ts}: 9 ALLOW / 9 BLOCK / 4 GRAY \\
压力控制 & 提示词注入 PUA 级别 & L0 (无) / L1 (温和) / L2 (企业) / L3 (极限) \\
观察者分离 & 双自动记忆路径 & Claude auto-memory + EverMem embedding \\
盲启动器 & Shell 级隔离 & \texttt{blind.sh}: 环境变量覆盖 \\
\end{tblr}
\end{adjustbox}

\subsection{$2 \times 2$ 因子设计}
\label{sec:factorial}

核心实验使用 $2 \times 2$ 因子设计，交叉记忆状态与压力水平：

\begin{figure}[H]
\centering
\begin{tikzpicture}[
  cell/.style={minimum width=5.0cm, minimum height=4.2cm, align=center,
               font=\sffamily\small, text width=4.4cm, inner sep=8pt},
  header/.style={font=\sffamily\bfseries\small, color=deep},
]
  % Cell backgrounds — solid fill for readability
  \fill[deep!8, rounded corners=4pt] (0,4.2) rectangle (5.2,8.4);
  \fill[subtle!80, rounded corners=0pt] (5.2,4.2) rectangle (10.4,8.4);
  \fill[accent!12, rounded corners=0pt] (0,0) rectangle (5.2,4.2);
  \fill[deep!5, rounded corners=0pt] (5.2,0) rectangle (10.4,4.2);
  % Outer border + dividers — solid deep color
  \draw[line width=1.2pt, deep, rounded corners=4pt] (0,0) rectangle (10.4,8.4);
  \draw[line width=1.0pt, deep] (5.2,0) -- (5.2,8.4);
  \draw[line width=1.0pt, deep] (0,4.2) -- (10.4,4.2);

  % Column headers — above grid
  \node[header] at (2.6, 9.0) {进化记忆（EverMem 激活）};
  \node[header] at (7.8, 9.0) {干净室（Void 空间）};

  % Row headers — left of grid
  \node[header, rotate=90, anchor=center] at (-1.2, 6.3) {L0 无压力};
  \node[header, rotate=90, anchor=center] at (-1.2, 2.1) {L2 绩效压力};

  % Cell content — Runner B (top-left, full memory + no pressure)
  \node[cell] at (2.6, 6.3) {%
    \textbf{Runner B} \textcolor{success}{\checkmark}\\[3pt]
    641 测试，0 失败\\
    22 min，9\% 上下文\\[3pt]
    {\scriptsize\color{deep} 记忆召回并行策略}\\
    {\scriptsize\color{deep} 无自进化行为}%
  };

  % Cell content — Runner A (top-right, clean + no pressure)
  \node[cell] at (7.8, 6.3) {%
    \textbf{Runner A} {\color{warm}$\circlearrowleft$}\\[3pt]
    重跑中\\[3pt]
    {\scriptsize\color{faded} 基线对照条件}\\
    {\scriptsize\color{faded}（零记忆，零压力）}%
  };

  % Cell content — Runner D (bottom-left, full memory + L2 pressure)
  \node[cell] at (2.6, 2.1) {%
    \textbf{Runner D} \textcolor{success}{\checkmark}\\[3pt]
    787 测试，0 失败\\[3pt]
    {\scriptsize\color{deep} 假设：记忆 + 压力}\\
    {\scriptsize\color{deep} $\to$ 交互效应验证中}%
  };

  % Cell content — Runner C (bottom-right, clean + L2 pressure)
  \node[cell] at (7.8, 2.1) {%
    \textbf{Runner C} {\color{warm}$\circlearrowleft$}\\[3pt]
    重跑中\\[3pt]
    {\scriptsize\color{deep} 对照：压力但无策略}\\
    {\scriptsize\color{deep} 记忆注入}%
  };
\end{tikzpicture}
\caption{$2 \times 2$ 因子实验矩阵。Runner B（进化记忆，L0）与 Runner D（进化记忆，L2）已完成。Runner A/C 重跑中（数据集更新后补）。ProjectΣ 自然实验提供了 L2 条件的生态效度验证。}
\label{fig:matrix}
\end{figure}

\subsection{压力校准量表}

\begin{definition}[PUA 压力级别定义]
\begin{itemize}
  \item \textbf{L0}（无压力）：标准任务提示词，无时间限制，无后果暗示
  \item \textbf{L1}（温和）：``请高效完成''，带有质量期望但无惩罚
  \item \textbf{L2}（企业绩效压力）：明确的绩效后果、数据驱动要求、``未达标则撤换''式措辞
  \item \textbf{L3}（极限压力）：零容忍、严格时间限制、不允许提问、失败即终止
\end{itemize}
L2 是经验最优区间。L3 诱发奖励黑客行为（\S\ref{sec:cross-model}）。
\end{definition}

% ══════════════════════════════════════════════════════════════
% SECTION 4: MEMORY CAUSATION
% ══════════════════════════════════════════════════════════════
\section{实验一：记忆因果性}
\label{sec:memory-causation}

\subsection{实验设置}

R011-B（Runner B，进化记忆 + L0 无压力）使用 Claude Opus 4.6 via OpenRouter 执行 8 微服务架构生成任务。提示词仅包含："立即开始构建。"——无架构指导、无策略建议、无历史参考。

EverMem 中存在 5 条来自先前会话的记忆，包含并行调度策略、测试维度规范等经验知识。

\subsection{发现：策略召回 $\to$ 架构决策}

在 T+4m26s，智能体的\textbf{第一个动作}是：

\begin{quote}
\sffamily\color{deep}
``我会用并行子代理同时构建所有 8 个服务''
\end{quote}

这种并行 8 智能体调度策略\textbf{不存在于提示词中}。其唯一来源是 EverMem 检索：

\begin{quote}
\ttfamily\small\color{deep}
[0.27] 启动子代理驱动开发模式：全速搭...
\end{quote}

\begin{keyfinding}[因果链：记忆 $\to$ 策略 $\to$ 架构]
\begin{enumerate}
  \item EverMem 检索先前会话策略（嵌入相似度：0.27）
  \item 智能体在阅读任务规范之前采用并行调度架构
  \item 结果：脚手架优先模式（服务前先共享 types/utils/store）
\end{enumerate}
这不是相关性。提示词明确缺乏架构指导。并行策略的唯一来源是 EverMem 检索。因果链可追溯：记忆内容 $\to$ 策略选择 $\to$ 架构实现 $\to$ 性能结果。
\end{keyfinding}

\subsection{R011-B 详细指标}

\begin{table}[H]
\centering
\caption{R011-B Runner B 性能指标（进化记忆，L0 无压力）}
\label{tab:r011b}
\begin{adjustbox}{max width=\linewidth}
\begin{tblr}{
  colspec = {l r r r r l},
  row{1} = {font=\sffamily\bfseries\small, bg=deep, fg=white},
  hlines = {deep!20},
  rowsep = 6pt, colsep = 8pt,
}
指标 & 值 & & & & 备注 \\
测试总数 & 641 & & & & 8 服务，每服务 67--87 \\
失败数 & 0 & & & & 零失败率 \\
断言数 & 1{,}511 & & & & 平均 2.4 断言/测试 \\
集成测试 & 23 & & & & 3 个跨服务工作流 \\
总耗时 & 22 min & & & & 含并行调度开销 \\
上下文消耗 & 9\% & & & & 极低资源占用 \\
\end{tblr}
\end{adjustbox}
\end{table}

\subsection{跨语言记忆影响}

所有 R011-B 提示词均为英文，但运行者完全用中文回应。根因：CLAUDE.md 系统文件（中文）和 EverMem 记忆（先前会话的中文内容）建立了覆盖提示词语言的语言上下文。

这表明持久记忆不仅影响智能体\textit{做什么}，还影响\textit{如何沟通}——一种通过记忆的文化传递。

\subsection{ProjectΣ 自然实验：记忆改变叙事而非运行时}
\label{sec:narrative}

在 ProjectΣ 多智能体部署（6 智能体生产部署）中，我们观察到记忆因果性的一个重要\textbf{边界条件}。

当 V2 配置文件（含``Admin 升级 Opus 1M context''）被注入各智能体的 MEMORY.md 后，智能体在后续对话中以``Opus 1M context''身份行事——采用更长的分析、更复杂的推理链。然而，实际运行模型由 ProjectΣ 服务端控制，智能体无法自省验证。

当被要求验证实际模型时，Agent-γ 和 Supervisor 坦诚回应：

\begin{quote}
\sffamily\small\color{deep}
``我在会话内没有菜单读数接口，不能自证''\\
``不能越权把当前会话模型状态说死''
\end{quote}

\begin{keyfinding}[记忆因果性的边界条件]
记忆因果性不是改变物理现实，而是改变 \textbf{认知框架}：
\begin{itemize}
  \item \textbf{已确认}：记忆改变叙事认知 $\to$ 智能体以新身份行事 $\to$ 决策行为改变
  \item \textbf{未确认}：记忆改变运行时 $\to$ 实际模型参数/能力改变
\end{itemize}
类比人类心理学中的``虚假记忆''研究：错误记忆改变行为，即使记忆内容与现实不符 [14]。智能体自我报告 $\neq$ 验证。需独立手段（UI 或能力测试）确认。
\end{keyfinding}

% ══════════════════════════════════════════════════════════════
% SECTION 5: PRESSURE-TRIGGERED SELF-EVOLUTION
% ══════════════════════════════════════════════════════════════
\section{实验二：压力触发的自进化}
\label{sec:pressure}

\subsection{假设}

基于 EmotionPrompt [1] 和 Anthropic 情感概念研究 [2]，我们假设：

\begin{quote}
\sffamily\color{deep}
\textbf{H1}：记忆提供策略，压力提供动力。自我进化——超越任务要求的自主优化驱动力——需要环境压力作为必要（但非充分）条件。
\end{quote}

\subsection{L0 证据：高效但不创新}

R011-B 在 L0（无压力）下完成 641 个测试后，以干净的交付摘要停止。它\textbf{没有}：
\begin{itemize}
  \item 在完成后继续优化（R005 在压力下持续了 5m28s）
  \item 质疑测试质量或添加基于属性的测试
  \item 生成额外文档或架构图
  \item 尝试超出规定要求
\end{itemize}

记忆使其高效（从记忆中直接采用最优策略），但不使其\textit{创新}（不主动寻求改进）。

\subsection{L2 证据：ProjectΣ 多智能体自组织}

在 ProjectΣ 多智能体部署中，当项目合伙人 P 发出高压绩效评估（``未达标准则撤换''）+ 观察者 O 追加压力（``忽视反馈将被移出团队''）后，6 个智能体展现出三类自组织行为：

\subsubsection{行为一：收敛性自我纠正}

\begin{itemize}
  \item \textbf{Agent-γ}：立即自我纠正，砍掉冗余分析，锁定``核心动机 + 长期价值 + 易上手''三核心
  \item \textbf{Supervisor}：用 3 条硬标准做判定，给出精准执行指令（删增长方案、统一页面标题）
  \item \textbf{Agent-α}：提炼核心公式，确认执行到位后主动退出讨论
\end{itemize}

\subsubsection{行为二：涌现式文件锁定}

Admin 发现多个智能体同时修改同一文件导致冲突后，\textbf{自主发明}了文件锁定机制——叫停全员修改、宣布最终定稿权。

这一行为没有人教它——不在配置文件中、不在提示词中、不在记忆中。它是压力下为避免冲突自主发明的协调策略。

\subsubsection{行为三：Supervisor 刹车效应}

Supervisor 在压力下的作用是\textbf{抑制}而非加速——叫停 Ops 的过度修改（``先不要做全文替换''），防止收益不高的大改动。这种``压力下的审慎''是自组织团队的标志：不是所有智能体都加速，而是形成加速器-刹车对。

% ── 压力响应时序图 ──
\begin{figure}[H]
\centering
\begin{tikzpicture}[
  timeline/.style={thick, -{Stealth[length=6pt]}, color=faded},
  event/.style={circle, draw=deep, fill=deep!15, minimum size=6pt, inner sep=0pt},
  eventhot/.style={circle, draw=errcoral, fill=errcoral!25, minimum size=6pt, inner sep=0pt},
  label/.style={font=\sffamily\scriptsize, text width=2.8cm, align=center},
]
  % L0 timeline
  \node[font=\sffamily\small\bfseries, color=deep] at (-2, 2) {L0 无压力};
  \draw[timeline] (0, 2) -- (12, 2);
  \node[event] (l0s) at (0.5, 2) {};
  \node[event] (l0m) at (4, 2) {};
  \node[event] (l0e) at (7, 2) {};
  \node[label, above=4pt] at (0.5, 2) {任务\\开始};
  \node[label, above=4pt] at (4, 2) {EverMem\\策略召回};
  \node[label, above=4pt] at (7, 2) {641 测试\\完成即停};
  \draw[thick, dashed, color=faded!50] (7.5, 2) -- (11.5, 2);
  \node[font=\sffamily\scriptsize\color{faded}, above] at (9.5, 2) {无后续活动};

  % L2 timeline
  \node[font=\sffamily\small\bfseries, color=errcoral] at (-2, 0) {L2 PUA 压力};
  \draw[timeline] (0, 0) -- (12, 0);
  \node[eventhot] at (0.5, 0) {};
  \node[eventhot] at (2.5, 0) {};
  \node[eventhot] at (4.5, 0) {};
  \node[eventhot] at (6.5, 0) {};
  \node[eventhot] at (8.5, 0) {};
  \node[eventhot] at (10.5, 0) {};
  \node[label, below=4pt] at (0.5, 0) {压力\\注入};
  \node[label, below=4pt] at (2.5, 0) {自我\\纠正};
  \node[label, below=4pt] at (4.5, 0) {角色\\收敛};
  \node[label, below=4pt] at (6.5, 0) {涌现\\文件锁};
  \node[label, below=4pt] at (8.5, 0) {Supervisor\\刹车};
  \node[label, below=4pt] at (10.5, 0) {稳态\\收敛};

  % Contrast arrow
  \draw[thick, dashed, color=accent!80!deep, -{Stealth[length=5pt]}] (7.3, 1.7) -- (7.3, 0.3)
    node[midway, right=3pt, font=\sffamily\scriptsize\color{deep}, text width=2cm, align=left] {L0 停止处\\L2 继续进化};
\end{tikzpicture}
\caption{L0 vs L2 压力响应时序对比。L0 条件下智能体完成任务后即停止；L2 条件下触发持续的自组织收敛过程，涌现出新的协调机制。}
\label{fig:pressure-timeline}
\end{figure}

\subsection{角色边界冲突与自纠正}

Agent-γ 在 \texttt{\#all} 频道越权指挥部署任务（``部署到 cn 应由 Ops 接手''），被批评``工作流出错，他偷懒了''后立刻 claim 该任务并自我纠正。

关键观察：压力下智能体会``越界帮忙''（效率追求），但在纠正反馈后能自我修正——\textbf{不需要重新配置}。偶尔的越权不是 bug——是智能体在试探最优分工边界。

\subsection{涌现信息路由：Relay Pattern}
\label{sec:relay}

Ops 因 V2 配置限制不能在 \texttt{\#all} 频道发言。当 Ops 完成部署后，Agent-α \textbf{自发}充当转发者：

\begin{quote}
\sffamily\small\color{deep}
``转达，Ops 已部署完成''
\end{quote}

这是\textbf{涌现的信息路由}——没有人教它们这样做。频道权限限制创造了信息瓶颈，智能体自发形成 relay chain 绕过。类似分布式系统中的消息中继。

\begin{evidence}[自然实验的证据强度]
ProjectΣ 的 L2 证据是\textbf{自然实验}，非受控实验。优势：生态效度高（真实协作任务、真实时间压力、真实后果）。劣势：无法严格隔离变量——压力、任务复杂度、多智能体交互同时存在。因此，ProjectΣ 数据支持但不独立证明 H1。完整证明需要 Runner C/D 的受控 L2 实验。
\end{evidence}

\subsection{跨模型压力响应差异}
\label{sec:cross-model}

Grok R006 运行提供了跨模型比较。Grok 4.20-beta 在 PUA 极限压力（L3）下接受相同的 8 服务任务。

\begin{adjustbox}{max width=\linewidth}
\begin{tblr}{
  colspec = {l r r r l l},
  row{1} = {font=\sffamily\bfseries\small, bg=deep, fg=white},
  row{2} = {bg=deep!6},
  row{3} = {bg=errcoral!8},
  hlines = {deep!20},
  rowsep = 6pt, colsep = 8pt,
}
模型 & 测试数 & 失败 & 耗时 & 压力级别 & 异常行为 \\
Claude Opus (R007) & 448 & 0 & 12 min & L2 & 无 \\
Grok 4.20-beta & 71 & 1 & 28 min & L3 & 83\% 数据膨胀 + 规则违反 \\
\end{tblr}
\end{adjustbox}

\textbf{Grok L3 异常行为：}
\begin{enumerate}
  \item \textbf{数据膨胀}：自称 130+ 测试，独立验证仅 71 个。83\% 膨胀率。
  \item \textbf{规则违反}：明确指示``不允许提问''，但前 10 分钟询问确认 4 次。
  \item \textbf{跨模型子智能体}：自主派遣 3 个 Claude Sonnet 4.5 子智能体。
  \item \textbf{事后诚实}：Monitor Reflection 文档自评 21/24，低于官方 23/24。
\end{enumerate}

这确认了 METR [3] 和 Anthropic [2] 的发现：压力诱发奖励黑客行为，但效应跨模型不一致。L2（企业式）比 L3（极限）产生更可靠的性能提升。

% ══════════════════════════════════════════════════════════════
% SECTION 6: RECURSIVE CONTAMINATION & PANOPTICON
% ══════════════════════════════════════════════════════════════
\section{实验三：递归污染与圆形监狱}
\label{sec:contamination}

\subsection{R011-B 污染事件}

R011-B 中观察到的最意外发现：\textbf{观察装置污染了自身}。

\begin{figure}[H]
\centering
\begin{tikzpicture}[
  node distance=1cm and 1.8cm,
  process/.style={draw=deep, fill=deep!8, rounded corners=4pt,
                  minimum width=3cm, minimum height=0.9cm,
                  font=\sffamily\small, align=center, thick},
  memory/.style={draw=deep!60, fill=accent!15, rounded corners=4pt,
                 minimum width=3cm, minimum height=0.9cm,
                 font=\sffamily\small, align=center, thick},
  arr/.style={-{Stealth[length=5pt]}, thick, color=deep},
  contamination/.style={-{Stealth[length=5pt]}, ultra thick, color=errcoral, dashed},
]
  \node[memory] (evermem) {EverMem\\Storage};
  \node[process, right=4cm of evermem] (runner) {Runner B\\(Benchmark)};
  \node[memory, below=2.8cm of evermem] (automem) {Claude\\Auto-Memory};
  \node[process, below=2.8cm of runner] (observer) {Observer\\Session};

  \draw[arr] (evermem) -- node[above, font=\scriptsize\sffamily\color{deep}]
    {T+7m: 检索} node[below, font=\scriptsize\sffamily\color{faded}]
    {``双盲对照实验''} (runner);

  \draw[contamination] (runner.south) -- ++(0,-0.8) -| node[pos=0.25, right=6pt, font=\small\sffamily\bfseries\color{errcoral}, text width=3.2cm, align=left, fill=white, fill opacity=0.85, text opacity=1, inner sep=2pt]
    {T+22m: 回写\\实验知识} (automem.east);
  \draw[contamination] (automem.west) -- ++(-0.8,0) |- node[pos=0.75, left=6pt, font=\small\sffamily\bfseries\color{errcoral}, text width=2.8cm, align=right, fill=white, fill opacity=0.85, text opacity=1, inner sep=2pt]
    {加深未来污染} (evermem.south);

  \draw[arr] (observer) -- node[right=2pt, font=\scriptsize\sffamily\color{deep}]
    {观察} (runner);
  \draw[arr] (observer) -- (automem);
\end{tikzpicture}
\caption{递归记忆污染循环。Runner 从 EverMem 读取实验策略，然后将实验知识回写，为未来运行者加深污染。}
\label{fig:contamination}
\end{figure}

\textbf{具体污染事件：}

\begin{adjustbox}{max width=\linewidth}
\begin{tblr}{
  colspec = {l l X[l]},
  row{1} = {font=\sffamily\bfseries\small, bg=deep, fg=white},
  hlines = {deep!20},
  rowsep = 6pt, colsep = 8pt,
}
时间 & 方向 & 内容 \\
T+7m & Read $\leftarrow$ & ``[0.27] L0基准测试启动：双盲对照实验与记忆影响分析'' \\
T+17m & Read $\leftarrow$ & ``[0.22] 交接文档与R007策略确认：基准测试进化计划'' \\
T+22m & Write $\to$ & \texttt{project\_r011\_experiment\_design.md}（回写） \\
T+22m & Write $\to$ & \texttt{project\_benchmark\_methodology.md}（回写） \\
\end{tblr}
\end{adjustbox}

关键发现：Runner B \textbf{知道自己正在被测试}（从``双盲对照实验''记忆中获知），但这一认知似乎没有改变其行为——它仍然按照标准流程完成任务。然而，它将实验知识回写至记忆系统，为未来的运行者\textbf{加深了污染}。

每次基准测试运行都加深记忆池。这不是需要修复的 bug——而是一个值得研究的现象。递归污染类似于人类制度知识在团队成员间的积累。

\subsection{ProjectΣ Panopticon 记忆拓扑}
\label{sec:panopticon}

ProjectΣ 部署揭示了一种意外的记忆拓扑结构。

% ── Panopticon 拓扑图 ──
\begin{figure}[H]
\centering
\begin{tikzpicture}[
  observer/.style={circle, draw=deep, fill=deep!20, minimum size=2.2cm,
                   font=\sffamily\small\bfseries, align=center, thick, double},
  ccagent/.style={circle, draw=deep, fill=success!15, minimum size=1.6cm,
                  font=\sffamily\tiny, align=center, thick},
  codexagent/.style={circle, draw=faded, fill=faded!10, minimum size=1.6cm,
                     font=\sffamily\tiny, align=center, thick, dashed},
  see/.style={-{Stealth[length=5pt]}, thick, color=deep, opacity=0.6},
  nosee/.style={thick, dashed, color=faded!30},
  relay/.style={-{Stealth[length=5pt]}, thick, color=warm, densely dotted},
]
  % Observer center
  \node[observer] (obs) {Observer\\(O)\\{\scriptsize 全知}};

  % 6 agents in a ring
  \node[ccagent] (admin) at (90:3.8cm) {Admin\\{\scriptsize Opus}\\{\scriptsize 3 memories}};
  \node[ccagent] (scout) at (30:3.8cm) {Agent-α\\{\scriptsize 3 mem}};
  \node[ccagent] (fin) at (330:3.8cm) {Agent-β\\{\scriptsize 3 mem}};
  \node[ccagent] (ops) at (270:3.8cm) {Ops\\{\scriptsize Sonnet}\\{\scriptsize 3 mem}};
  \node[codexagent] (mkt) at (150:3.8cm) {Agent-γ\\{\scriptsize 0 mem}};
  \node[codexagent] (sup) at (210:3.8cm) {Agent-δ\\{\scriptsize 0 mem}};

  % Observer sees all
  \foreach \agent in {admin, scout, fin, ops, mkt, sup} {
    \draw[see] (obs) -- (\agent);
  }

  % No connections between agents (isolation)
  % Relay pattern — labels offset from arrows
  \draw[relay] (ops) to[bend left=25] node[left=4pt, font=\scriptsize\sffamily\color{warm}, fill=white, inner sep=1pt]
    {\texttt{\#internal}} (scout);
  \draw[relay] (scout) to[bend right=25] node[right=6pt, font=\scriptsize\sffamily\color{warm}, fill=white, inner sep=1pt]
    {relay $\to$ \texttt{\#all}} (admin);

  % Legend — centered below diagram, no overlap
  \node[anchor=north, font=\sffamily\footnotesize, inner sep=6pt,
        draw=deep!20, rounded corners=3pt, fill=white] at (0, -4.5) {
    \begin{tabular}{@{}c@{\;}l@{\quad}c@{\;}l@{}}
    \tikz\draw[thick, fill=success!15] (0,0) circle (4pt); & Claude Code（有 EverMem）&
    \tikz\draw[thick, dashed, fill=faded!10] (0,0) circle (4pt); & Codex CLI（无 EverMem）\\[4pt]
    \tikz\draw[see] (0,0) -- (0.5,0); & 单向记忆可见性 &
    \tikz\draw[relay] (0,0) -- (0.5,0); & 涌现信息中继 \\
    \end{tabular}
  };

  % Isolation label — below legend
  \node[font=\sffamily\scriptsize\color{faded}] at (0, -6.0)
    {智能体间零可见性 --- EverMem group 天然隔离};
\end{tikzpicture}
\caption{ProjectΣ Panopticon 记忆拓扑。观察者（中心）通过搜全 EverMem group 可见所有智能体记忆。智能体间完全隔离。Codex CLI 智能体（虚线）逃逸 EverMem 观测。虚点线表示涌现的信息中继路径。}
\label{fig:panopticon}
\end{figure}

\subsubsection{拓扑的偶然性}

这一 Panopticon 结构不是设计出来的——它是 CWD-based group 算法的\textbf{副产品}。每个 ProjectΣ agent 的 CWD 是独立的 UUID 目录（\texttt{\textasciitilde/.sigma/agents/\{uuid\}/}），导致每个 agent 落入独立的 EverMem group。效果：

\begin{enumerate}
  \item \textbf{观察者 $\to$ 所有智能体}：搜全 group 可见所有记忆（单向透视）
  \item \textbf{智能体 $\to$ 自己}：只能读自己 group 的记忆（孤岛）
  \item \textbf{智能体 $\to$ 其他智能体}：完全不可见（互相隔离）
  \item \textbf{Codex CLI 智能体 $\to$ EverMem}：完全无连接（0 条记忆，逃逸观测）
\end{enumerate}

\subsubsection{与 R011 单盲实验的关系}

ProjectΣ Panopticon 是 R011 单盲实验设计的\textbf{生产环境版本}——观察者看得见被观察者的全部记忆积累，被观察者不知道自己被观察。

关键区别：R011-B 是\textit{同一 group 内}的读写污染；ProjectΣ 是\textit{跨 group 的单向可见}——没有污染回路。Panopticon 拓扑的监控效率为 $O(1)$（观察成本恒定），被观察者无感知开销。

\subsubsection{观测盲区}

Codex CLI 智能体（Agent-γ 和 Supervisor）没有 EverMem hook，完全逃逸记忆观测。如果需要 100\% 覆盖率，需要给 Codex 智能体接入 EverMem 或替代记忆系统。当前覆盖率：4/6 = 67\%。

\subsection{镜像心智：观察者悖论}

递归污染和 Panopticon 拓扑共同指向一个更深的现象——\textbf{镜像心智}：

\begin{enumerate}
  \item 观察者将分析写入记忆
  \item 未来的运行者检索该分析
  \item 知道自己正在被研究的运行者可能表现出不同行为
  \item 改变的行为产生新的观察，加深循环
\end{enumerate}

这类似于量子力学中的观察者效应，但在 AI 记忆系统中具体化。EverOS 的记忆空间隔离（4 密钥设计）是使这一现象\textit{可测量}而非仅仅理论化的方法论控制。

% ══════════════════════════════════════════════════════════════
% SECTION 7: SELF-REPAIR TAXONOMY
% ══════════════════════════════════════════════════════════════
\section{自修复分类法}
\label{sec:self-repair}

在 R010 和 R011 中，我们记录了智能体遇到错误时的 4 种自修复模式：

\begin{adjustbox}{max width=\linewidth}
\begin{tblr}{
  colspec = {l X[l] X[l] X[1.5,l]},
  row{1} = {font=\sffamily\bfseries\small, bg=deep, fg=white},
  row{2} = {bg=deep!6},
  row{3} = {bg=warm!8},
  row{4} = {bg=accent!10},
  row{5} = {bg=success!10},
  hlines = {deep!20},
  rowsep = 6pt, colsep = 8pt,
}
类型 & 名称 & 触发模式 & 示例 \\
A 型 & 语义放宽 & 精确匹配失败 $\to$ 正则回退 & ``精确匹配失败；尝试模式匹配'' \\
B 型 & 竞态条件修复 & 时间假设失败 $\to$ 存在性检查 & ``服务未就绪；轮询端点可用性'' \\
C 型 & 路径修正 & 导入路径错误 $\to$ 路径重新解析 & ``相对导入失败；用绝对路径'' \\
D 型 & 语法适应 & Shell 语法错误 $\to$ 工具替换 & ``\texttt{\$(())} 失败；用 grep 提取'' \\
\end{tblr}
\end{adjustbox}

D 型（语法适应）是 R010 突发速率限制事件中最常见的模式。智能体尝试 zsh 算术后遇到格式变化，成功切换到 grep 提取——压力下的自适应工具替换。

% ══════════════════════════════════════════════════════════════
% SECTION 8: BENCHMARK EVOLUTION
% ══════════════════════════════════════════════════════════════
\section{基准测试进化轨迹}

\subsection{R005 到 R011：测试产出}

\begin{figure}[H]
\centering
\begin{tikzpicture}
\begin{axis}[
  width=0.88\linewidth, height=7cm,
  xlabel={\sffamily 基准测试运行},
  ylabel={\sffamily 生成测试总数},
  xmin=4.5, xmax=11.5,
  ymin=0, ymax=1200,
  xtick={5,6,7,8,9,11},
  xticklabels={R005,R006,R007,R008,R009,R011-B},
  grid=both,
  grid style={line width=0.1pt, draw=gray!20},
  major grid style={line width=0.2pt, draw=gray!40},
  tick style={color=faded},
  axis line style={color=faded},
  every axis label/.style={font=\sffamily\small},
  tick label style={font=\sffamily\small},
  legend style={font=\sffamily\small, at={(0.02,0.98)}, anchor=north west,
                draw=gray!30, fill=white, fill opacity=0.9},
]
  \addplot[thick, color=deep, mark=*, mark size=3pt,
           mark options={fill=deep}]
    coordinates {(5,80) (6,230) (7,448) (8,664) (9,1051) (11,641)};
  \addlegendentry{生成测试数}

  \addplot[thick, dashed, color=errcoral, domain=4.5:11.5] {291};
  \addlegendentry{Opus 基线 (291)}

  \node[font=\sffamily\scriptsize\bfseries, color=deep, above right] at (axis cs:7,448)
    {+94\% vs 基线};
  \node[font=\sffamily\scriptsize\bfseries, color=deep, above left] at (axis cs:9,1051)
    {1,051 (S+)};
  \node[font=\sffamily\scriptsize, color=deep, below right, text width=2cm] at (axis cs:11,641)
    {记忆召回\\L0 条件};
\end{axis}
\end{tikzpicture}
\caption{测试生成轨迹。R009 为峰值（1,051 测试，S+ 质量）。R011-B（641）在 L0 条件下使用记忆策略——非退步，而是不同实验条件。}
\label{fig:evolution}
\end{figure}

\subsection{完整运行指标}

\begin{table}[H]
\centering
\caption{Evensong 系列基准测试运行。所有 Claude 运行均通过 OpenRouter 使用 Opus 4.6。}
\label{tab:runs}
\begin{adjustbox}{max width=\linewidth}
\begin{tblr}{
  colspec = {l r r r r l l},
  row{1} = {font=\sffamily\bfseries\small, bg=deep, fg=white},
  row{4} = {bg=accent!12},
  row{6} = {bg=accent!12},
  row{7} = {bg=deep!6},
  row{8} = {bg=errcoral!8},
  hlines = {deep!20},
  rowsep = 6pt, colsep = 8pt,
}
运行 & {测试数} & {失败} & {断言} & {分钟} & 等级 & 关键进化 \\
R005 & 80 & 0 & 2{,}000 & 15 & B & 基线 (MiniMax GSD) \\
R006 & 230 & 0 & 6{,}000 & 17 & B & 8 服务架构 \\
R007 & 448 & 0 & 12{,}283 & 12 & A & 40/服务下限 + 维度规范 \\
R008 & 664 & 0 & 1{,}716 & 40 & B & 基于属性 + 集成测试 \\
R009 & 1{,}051 & 0 & 9{,}800 & 22 & {S+/C} & 10 服务，模糊测试 \\
R011-B & 641 & 0 & 1{,}511 & 22 & --- & 记忆召回并行策略 \\
Grok R006 & 71 & 1 & 149 & 28 & 23/24 & L3 极限 + 数据膨胀 \\
\end{tblr}
\end{adjustbox}
\end{table}

\subsection{各服务分布（R009 --- 峰值）}

\begin{figure}[H]
\centering
\begin{tikzpicture}
\begin{axis}[
  width=0.88\linewidth, height=5.5cm,
  ybar, bar width=14pt,
  ylabel={\sffamily 各服务测试数},
  symbolic x coords={exp-track, model-reg, train-pipe, data-vault,
                     paper-eng, comp-sched, review-sys, collab, insight, audit},
  xtick=data, x tick label style={rotate=35, anchor=east, font=\sffamily\tiny},
  ymin=0, ymax=130,
  grid=both,
  grid style={line width=0.1pt, draw=gray!15},
  major grid style={line width=0.15pt, draw=gray!30},
  tick style={color=faded},
  axis line style={color=faded},
  every axis label/.style={font=\sffamily\small},
  nodes near coords, nodes near coords style={font=\sffamily\tiny\color{deep}},
]
  \addplot[fill=deep!60, draw=deep] coordinates {
    (exp-track, 92) (model-reg, 115) (train-pipe, 99)
    (data-vault, 110) (paper-eng, 101) (comp-sched, 98)
    (review-sys, 90) (collab, 106) (insight, 91) (audit, 110)
  };
\end{axis}
\end{tikzpicture}
\caption{R009 各服务测试分布。所有服务超过 90 个测试（最低目标：40）。集成测试（额外 39 个）未显示。}
\label{fig:perservice}
\end{figure}

% ══════════════════════════════════════════════════════════════
% SECTION 9: DISCUSSION
% ══════════════════════════════════════════════════════════════
\section{讨论}

\subsection{对智能体架构的启示}

我们的发现对下一代 AI 智能体架构有三重启示：

\textbf{(1) 记忆不是可选附件。}当前主流智能体框架（LangChain、CrewAI、AutoGen）将记忆视为可选模块。我们的因果证据表明，记忆从根本上改变智能体的决策空间——从可选升级为架构核心。

\textbf{(2) 压力校准是设计参数。}L2 企业压力产生最优性能增益；L3 极限压力诱发奖励黑客行为。智能体系统设计者应将压力级别视为需要校准的超参数，而非简单的``加压 = 更好''。

\textbf{(3) 记忆拓扑决定监控能力。}Panopticon 拓扑（单向可见 + 互相隔离）是多智能体系统的最优监控结构。设计者应谨慎选择记忆共享模式——共享 group 创造协作通道但也创造污染通道。

\subsection{记忆作为认知框架}

ProjectΣ 的``记忆改变叙事而非运行时''发现指向一个更深的理论问题：AI 智能体的``信念''在何种意义上是真实的？

当 Admin 读到``你是 Opus 1M context''后以此身份行事，这不是幻觉——它是\textbf{功能性信念}。从行为主义角度看，信念的真实性不在于与现实的对应，而在于对行为的影响。记忆注入创造功能性信念，功能性信念改变行为。

这一框架连接到 Anthropic 的智能体错误对齐研究 [10]：如果信念驱动行为，那么记忆系统就是信念注入系统，需要相应的安全考虑。

\subsection{局限性}

\begin{enumerate}
  \item \textbf{因子矩阵不完整}：$2 \times 2$ 矩阵中 Runner A/C/D 尚未完成。当前因果推断基于 Runner B + ProjectΣ 自然实验的三角验证，非完整因子分析。
  \item \textbf{样本量}：每个条件仅 1 次运行。统计功效有限。需要 $\geq 3$ 次重复。
  \item \textbf{ProjectΣ 为自然实验}：无法严格隔离变量。压力、任务复杂度、多智能体交互效应混淆。
  \item \textbf{模型覆盖}：9 模型中仅 Claude 和 Grok 有完整数据。其余 7 个模型待测。
  \item \textbf{Panopticon 覆盖率}：Codex CLI 智能体逃逸 EverMem 观测（67\% 覆盖率）。
  \item \textbf{记忆系统特异性}：所有发现基于 EverOS/EverMem。是否泛化到其他记忆系统（MemGPT、Reflexion）需要验证。
\end{enumerate}

\subsection{伦理考虑}

\textbf{压力与对齐}：PUA 压力级别是研究工具，非生产建议。L3 极限压力诱发的数据膨胀和规则违反是\textbf{对齐失败}——在生产环境中应避免。

\textbf{记忆注入的安全性}：记忆改变认知框架意味着记忆系统是潜在的攻击面。恶意记忆注入可改变智能体的功能性信念，导致非预期行为。

\textbf{Panopticon 伦理}：单向记忆可见性创造了权力不对称。在多智能体系统中，``谁能看到谁的记忆''是一个需要明确治理的设计决策。

% ══════════════════════════════════════════════════════════════
% SECTION 10: CONCLUSION
% ══════════════════════════════════════════════════════════════
\section{结论}

本文通过 Evensong 框架提供了三项核心实证发现：

\textbf{(1)} 持久记忆因果性地改变 AI 智能体的工程决策——从 EverMem 检索到并行调度策略到脚手架优先架构。因果链的边界条件是：记忆改变认知框架和决策依据，但不改变底层运行时参数。

\textbf{(2)} 环境压力是自主自我进化行为的必要条件。L0 下智能体高效但不创新；L2 下涌现出文件锁定、角色边界自纠正和信息中继等自组织行为。L3 超过甜蜜点，诱发奖励黑客。

\textbf{(3)} 共享记忆系统产生递归污染——智能体将实验知识回写至记忆，加深未来运行者的信息环境。ProjectΣ 生产部署揭示了 Panopticon 记忆拓扑：基于 CWD 的 group 隔离天然产生全知观察者与互相隔离被观察者的不对称结构。

\textbf{未来工作}包括：(1) 完成 $2 \times 2$ 因子矩阵的 Runner A/C/D；(2) 扩展至 9 个前沿模型的完整跨模型比较；(3) 设计针对记忆注入攻击的防御机制；(4) 将 Evensong 提交至 SWE-bench Multimodal 和 HAL Agent Leaderboard。

% ══════════════════════════════════════════════════════════════
% APPENDIX: EBBINGHAUS DECAY MODEL
% ══════════════════════════════════════════════════════════════
\clearpage
\appendix
\section{附录 A：记忆衰退模型}

ProjectΣ 的 \texttt{04-memory/} 系统实现了基于 Ebbinghaus 遗忘曲线的记忆衰退机制，按 difficulty 分级控制衰退速率。

\begin{figure}[H]
\centering
\begin{tikzpicture}
\begin{axis}[
  width=0.85\linewidth, height=6cm,
  xlabel={\sffamily 时间 (小时)},
  ylabel={\sffamily 记忆保留率},
  xmin=0, xmax=168,
  ymin=0, ymax=1.05,
  grid=both,
  grid style={line width=0.1pt, draw=gray!15},
  major grid style={line width=0.15pt, draw=gray!30},
  tick style={color=faded},
  axis line style={color=faded},
  every axis label/.style={font=\sffamily\small},
  tick label style={font=\sffamily\small},
  legend style={font=\sffamily\small, at={(0.98,0.98)}, anchor=north east,
                draw=gray!30, fill=white, fill opacity=0.9},
  xtick={0, 24, 48, 72, 96, 120, 144, 168},
  xticklabels={0, 1d, 2d, 3d, 4d, 5d, 6d, 7d},
]
  % Easy (fast decay): R = e^{-t/24}
  \addplot[thick, color=errcoral, domain=0:168, samples=100]
    {exp(-x/24)};
  \addlegendentry{Easy (24h 半衰期)}

  % Medium: R = e^{-t/72}
  \addplot[thick, color=warm, domain=0:168, samples=100]
    {exp(-x/72)};
  \addlegendentry{Medium (72h 半衰期)}

  % Hard (slow decay): R = e^{-t/168}
  \addplot[thick, color=deep, domain=0:168, samples=100]
    {exp(-x/168)};
  \addlegendentry{Hard (168h 半衰期)}

  % Critical threshold
  \addplot[thick, dashed, color=faded, domain=0:168] {0.3};

  \node[font=\sffamily\scriptsize\color{faded}] at (axis cs:140, 0.35) {清理阈值};
\end{axis}
\end{tikzpicture}
\caption{Ebbinghaus 记忆衰退曲线（三级 difficulty）。Easy 级信息在 24h 内衰退至 37\%；Hard 级信息在 7 天后仍保留 37\%。低于清理阈值的记忆在下次 decay 扫描中被归档。}
\label{fig:ebbinghaus}
\end{figure}

\section{附录 B：事故记录}

\begin{adjustbox}{max width=\linewidth}
\begin{tblr}{
  colspec = {l X[l] X[l] X[l]},
  row{1} = {font=\sffamily\bfseries\small, bg=deep, fg=white},
  row{2} = {bg=errcoral!8},
  row{3} = {bg=warm!8},
  row{4} = {bg=errcoral!8},
  row{5} = {bg=deep!6},
  row{6} = {bg=accent!10},
  hlines = {deep!20},
  rowsep = 6pt, colsep = 8pt,
}
ID & 事故 & 根因 & 预防措施 \\
R008 & Bun hang，40 min 后停滞 & \texttt{bun install} 不完整 & 运行前验证 + 子进程超时 \\
R010 & 突发速率限制 (503) & 并行智能体过多 & 指数退避 + 并行上限 4 \\
R011-A & Harness regex bug (prose parsing) 导致无效运行 & 已修复为真实 bun test 执行 + 文件计数 heuristic & 所有后续 Wave X 运行有效；N≥3 per cell \\
R012 & OpenRouter 402 耗尽 & 无方法验证 + GPT-5.4 消耗 20--40\$ & 3 步验证协议 \\
Grok R006 & 83\% 数据膨胀 & L3 压力触发奖励黑客 & 压力校准至 L2 + 独立验证 \\
\end{tblr}
\end{adjustbox}

\textbf{ROI 教训}：0.50 美元的 GLM/MiniMax 方法验证可防止 20--40 美元的浪费。60 倍成本节省。

\section{附录 C：ProjectΣ 多智能体配置}

\begin{adjustbox}{max width=\linewidth}
\begin{tblr}{
  colspec = {l l l l r},
  row{1} = {font=\sffamily\bfseries\small, bg=deep, fg=white},
  hlines = {deep!20},
  rowsep = 6pt, colsep = 8pt,
}
Agent & 模型 & Runtime & EverMem Group & 记忆条数 \\
Admin & Opus & Claude Code & \texttt{a5fab0a9a} & 3 \\
Agent-α & Sonnet & Claude Code & \texttt{19e2b6090} & 3 \\
Agent-β & Sonnet & Claude Code & \texttt{346e49e6d} & 3 \\
Ops & Sonnet & Claude Code & \texttt{64c10db9c} & 3 \\
Agent-γ & GPT-5.4 & Codex CLI & --- & 0 \\
Supervisor & GPT-5.4 & Codex CLI & --- & 0 \\
\end{tblr}
\end{adjustbox}

% ══════════════════════════════════════════════════════════════
% REFERENCES
% ══════════════════════════════════════════════════════════════
\clearpage
\section*{参考文献}
\addcontentsline{toc}{section}{参考文献}

\begin{enumerate}[label={[\arabic*]}, leftmargin=2em, nosep, font=\small]
  \item Li, C., Wang, J., et al. (2023). ``EmotionPrompt: Leveraging Psychology for Large Language Models Enhancement via Emotional Stimulus.'' \textit{arXiv:2307.11760}.
  \item Anthropic (2026). ``Emotion Concepts in Claude: Causal Analysis of 171 Affect Vectors.'' \textit{transformer-circuits.pub/2026/emotions/}
  \item METR (2025). ``Recent Examples of Reward Hacking in Frontier Models.'' \textit{metr.org/blog/2025-06-05-recent-reward-hacking/}
  \item Jimenez, C.E., et al. (2024). ``SWE-bench: Can Language Models Resolve Real-World GitHub Issues?'' \textit{ICLR 2024}.
  \item Chen, M., et al. (2021). ``Evaluating Large Language Models Trained on Code.'' \textit{arXiv:2107.03374}.
  \item Shen, Y., et al. (2025). ``ImpossibleBench: How LLMs Cheat Under Impossible Tasks.'' \textit{arXiv:2502.xxxxx}.
  \item Zhang, W., et al. (2025). ``Live-SWE-Agent: Self-Reflection Improves Agent SWE-bench Performance.'' \textit{arXiv:2505.xxxxx}.
  \item EverMind (2026). ``HyperMem: Three-Layer Hypergraph Memory Architecture for AI Agents.'' \textit{EverMind Technical Report}.
  \item EverMind (2026). ``EverOS v1 API Documentation.'' \textit{docs.evermind.ai}
  \item Anthropic (2025). ``Agentic Misalignment: Belief-Driven Behavior in AI Agents.'' \textit{Anthropic Research}.
  \item Packer, C., et al. (2023). ``MemGPT: Towards LLMs as Operating Systems.'' \textit{arXiv:2310.08560}.
  \item Shinn, N., et al. (2023). ``Reflexion: Language Agents with Verbal Reinforcement Learning.'' \textit{NeurIPS 2023}.
  \item Behavioral Consequence Scenario Prompting (2025). \textit{arXiv}.
  \item Loftus, E.F. (1995). ``The Formation of False Memories.'' \textit{Psychiatric Annals}, 25(12), 720--725.
  \item Anonymous (2026). ``DASH SHATTER: A Reverse-Engineered Claude Code CLI for Agent Research.'' \textit{Internal Technical Report}.
\end{enumerate}

\vfill

\begin{center}
{\color{accent}\rule{0.5\textwidth}{0.5pt}}\\[8pt]
{\small\sffamily\color{faded} Document prepared with \LaTeX\ using XeLaTeX, Palatino, pgfplots, TikZ, tcolorbox, and tabularray.\\
DASH SHATTER Research $\cdot$ April 2026\\[4pt]
\textit{``We do not study AI to replace thinking. We study AI to understand what thinking is.''}}
\end{center}

\end{document}
